\chapter{FT01 to FT05: what each fine-tuning run taught}\label{ch:ftruns}
\begingroup\sloppy % long provenance tags need loose lines rather than overfull ones

\begin{onething}
Six fine-tuning runs (FT01, FT02, FT02b, FT02c, FT03, FT05) and the release candidate FT04 bought better masks,
native bands and better reconstruction on all three sensors - and no measurable gain in detection by the model
alone. Each run was judged by criteria written before launch, and each taught one lesson: a frozen encoder's
ceiling, a gate read at the wrong checkpoint, edges that are also targets, a target that hides the anomaly, and a
score - residual Mahalanobis - that absorbs whatever better reconstruction changes.
\end{onething}

\section{Where the runs start}\label{sec:ftruns:start}

\Cref{ch:finetune} laid out the curriculum: a \emph{probe} (train only a small head on a frozen encoder), an
\emph{adapt} stage (continue masked reconstruction on the real sensors), a \emph{joint} stage (reconstruction and
segmentation together), optional pseudo-labels, and finally product use of the masks. It also taught the two
mechanical fixes that the early runs forced - the BatchNorm see-saw and mixed batches. This chapter reads the
runs themselves, in the order they ran on an L40S GPU (\texttt{g6e.2xlarge}) between 2026-10-16 and 2026-10-18,
as a sequence of lessons. \Cref{tab:ftruns:runs} is the map.

\begin{table}[htbp]
\centering
\caption{The fine-tuning runs of 2026-10-16 to 18. ``Verdict'' is against the criteria each run's README fixed
before launch.}
\label{tab:ftruns:runs}
\small
\begin{tabular}{llp{5.0cm}cl}
\toprule
run & stage & what changed & hours & verdict \\
\midrule
FT01 & probe & segmentation head on frozen F001 (epoch 109), EnMAP & 1.4 & passed (water 0.82) \\
FT02 & adapt & reconstruction on three sensors, one sensor per batch & 4.25 & gate not met \\
FT02b & adapt & FT02 continued with mixed-sensor batches & 2.2 & gate not met \\
FT02c & adapt & FT02b continued with an edge term & 2.2 & not met \\
FT03 & joint & reconstruction + segmentation, mixed batches & 3.0 & passed \\
FT05 & joint & implants in the input, clean tile as target & -- & not passed \\
FT04 & joint & pseudo-labels on PRISMA / AVIRIS-NG, rare-class weights & -- & passed \\
\bottomrule
\end{tabular}

\smallskip
{\footnotesize Source: \tagM{research/experiments/indradhanu-finetune/README.md} and the run READMEs beside it.}
\end{table}

\paragraph{Why F001 stopped at epoch 113.} \Cref{ch:pretrain} described F001 as a 200-epoch run. It did not
finish. The PRISMA bench under residual Mahalanobis read 0.861 at epoch 60 and 0.862 at epoch
109\tagM{research/experiments/indradhanu-se-pretrain/F001-final-pretrain/README.md}. Forty-nine epochs for 0.001:
the run was stopped at epoch 113 by decision, and every fine-tuning run starts, directly or through its parent,
from \file{F001_ep109.pt}.

\paragraph{Gates written before launch.} \index{pass criterion} Every run's README states, before launch, what
it expects and what would count as passing. That habit is the spine of this chapter. A criterion written after
seeing the numbers can always be made to pass; a criterion written before is a prediction that can be wrong, and
several of the lessons below are exactly a prediction being wrong.

\begin{analogy}
A pianist moves from an upright to a concert grand. After a week of practice the scales are cleaner on every
octave - that is reconstruction. In the recital, the only thing the audience hears, one piece is very slightly
worse - that is detection. Cleaner scales are evidence of progress, not a guarantee of a better recital; you still
have to play the recital. Every run below improved its scales. Very few changed the recital.
\end{analogy}

\section{FT01: the probe measures what the encoder already knows}\label{sec:ftruns:ft01}

\index{probe}
A \emph{probe} trains only the small segmentation head (\cref{ch:seghead}) on top of a frozen encoder - frozen
meaning its parameters, BatchNorm statistics included, do not change. Whatever the head learns must already be
present in the frozen features, so the probe measures what the encoder knows. FT01 ran 15 epochs on EnMAP in
1.4 hours\tagM{research/experiments/indradhanu-finetune/FT01-probe/README.md}.

\paragraph{Counting the head by hand.} The FT01 README calls the head ``about 372k parameters''. Predict the
exact figure. The head is a SegFormer-style decoder (\cref{ch:decoder}) of width 128 over the encoder's four
stages, whose channel counts are 32, 64, 160 and 256. Its layers, read from the shipped checkpoint's shapes:

\begin{worked}[label={we:ftruns:seghead}]{The segmentation head's parameters}
Recall from \cref{ch:shapes} that a convolution with $C_\text{in}$ inputs, $C_\text{out}$ outputs and a $k
\times k$ kernel has $C_\text{out} \cdot C_\text{in} \cdot k^2$ weights and $C_\text{out}$ biases. Do each line
on paper.
\begin{enumerate}
\item Four $1 \times 1$ projections to 128 channels. Weights:
  \[
  128 \times (32 + 64 + 160 + 256) = 128 \times 512 = 65{,}536 .
  \]
  Biases: $4 \times 128 = 512$. Total $65{,}536 + 512 = 66{,}048$.
\item A $1 \times 1$ fuse convolution from the four stacked projections, $4 \times 128 = 512$ channels, to 128:
  \[
  128 \times 512 + 128 = 65{,}536 + 128 = 65{,}664 .
  \]
\item A $3 \times 3$ refine convolution, 128 to 128:
  \[
  128 \times 128 \times 9 + 128 = 147{,}456 + 128 = 147{,}584 .
  \]
\item A $3 \times 3$ output convolution, 128 to 80. Why 80? Five classes, each needing a $4 \times 4$ block of
  pixels per quarter-resolution position for the pixel shuffle (\cref{ch:decoder}): $5 \times 16 = 80$.
  \[
  80 \times 128 \times 9 + 80 = 92{,}160 + 80 = 92{,}240 .
  \]
\end{enumerate}
Sum:
\begin{align*}
66{,}048 + 65{,}664 &= 131{,}712\\
131{,}712 + 147{,}584 &= 279{,}296\\
279{,}296 + 92{,}240 &= 371{,}536
\end{align*}
The shipped checkpoint's \texttt{seg\_head.*} tensors hold exactly 371,536 numbers
(\file{research/textbook/tools/compute_ftruns.py})\tagC. ``About 372k'' is right.
\end{worked}

\paragraph{The measure.} \index{intersection over union} The head is judged by IoU, intersection over union
(\cref{ch:seghead}): for one class, the number of pixels both predicted and labelled, divided by the number
predicted or labelled. If 6 pixels are labelled water, 5 are predicted water, and 4 are both,
\[
\text{IoU} = \frac{4}{6 + 5 - 4} = \frac{4}{7} = 0.571 .
\]
(The union is $6 + 5 - 4$: adding the two sets counts the overlap twice, so subtract it once.)

\begin{table}[htbp]
\centering
\caption{FT01: per-class IoU on EnMAP's 950 test tiles. best.pt is epoch 14 (lowest segmentation validation loss).}
\label{tab:ftruns:ft01}
\small
\begin{tabular}{lcccccc}
\toprule
epoch & water & cloud & thin cloud & shadow & snow & mean \\
\midrule
1 & 0.705 & 0.600 & 0.067 & 0.000 & 0.000 & 0.274 \\
3 & 0.777 & 0.733 & 0.076 & 0.000 & 0.000 & 0.317 \\
6 & 0.780 & 0.784 & 0.065 & 0.104 & 0.000 & 0.347 \\
9 & 0.815 & 0.766 & 0.061 & 0.185 & 0.000 & 0.365 \\
12 & 0.822 & 0.807 & 0.089 & 0.165 & 0.055 & 0.388 \\
\textbf{14} & \textbf{0.822} & \textbf{0.801} & \textbf{0.091} & \textbf{0.199} & 0.175 & \textbf{0.418} \\
15 & 0.823 & 0.800 & 0.089 & 0.193 & 0.169 & 0.415 \\
\bottomrule
\end{tabular}

\smallskip
{\footnotesize Source: \tagM{research/experiments/indradhanu-finetune/FT01-probe/README.md} (from \texttt{finetune\_metrics.csv}).}
\end{table}

\begin{figure}[htbp]
\centering
\begin{tikzpicture}
\begin{axis}[width=0.85\linewidth, height=6cm, xlabel={epoch}, ylabel={IoU on EnMAP test},
  xmin=0, xmax=16, ymin=0, ymax=0.9, xtick={1,3,6,9,12,14,15},
  legend style={at={(0.98,0.55)},anchor=east,font=\small}, grid=major, grid style={gray!20}]
\addplot[spec, thick, mark=*] coordinates {(1,0.705) (3,0.777) (6,0.780) (9,0.815) (12,0.822) (14,0.822) (15,0.823)};
\addplot[spat, thick, mark=square*] coordinates {(1,0.600) (3,0.733) (6,0.784) (9,0.766) (12,0.807) (14,0.801) (15,0.800)};
\addplot[loss, thick, mark=triangle*] coordinates {(1,0.000) (3,0.000) (6,0.104) (9,0.185) (12,0.165) (14,0.199) (15,0.193)};
\addplot[maskgrey, very thick, mark=diamond*] coordinates {(1,0.067) (3,0.076) (6,0.065) (9,0.061) (12,0.089) (14,0.091) (15,0.089)};
\legend{water, cloud, shadow, thin cloud}
\end{axis}
\end{tikzpicture}
\caption{FT01's IoU by epoch, drawn to scale from \cref{tab:ftruns:ft01}. Water, cloud and shadow climb; thin cloud
(grey) stays between 0.06 and 0.09 for the whole run.}
\label{fig:ftruns:ft01}
\end{figure}

The gate was ``water at least 0.6, or the plan is revisited''. Water reached 0.82 and cloud 0.80 (the README had
expected cloud at least 0.5). The interesting row is thin cloud (\cref{fig:ftruns:ft01}). Every other class
climbed over the 15 epochs; thin cloud sat at 0.06 to 0.09 from start to end, against an expectation of 0.3 to
0.5. The label audit of \cref{ch:seghead} had already found that EnMAP's cirrus labels are often driven by the
surface, not by cloud.

\begin{keypoint}
A head that cannot improve on one class while it improves on every other is usually hitting the labels' ceiling,
not the encoder's: the encoder and head are shared by every class, so a feature limit would still let the class
move a little.
\end{keypoint}

So the FT01 README says: do not tune later runs on thin cloud. Snow's 0.17 rests on a single test tile and is not a
measurement.

\section{FT02: better reconstruction, a gate missed by 0.006}\label{sec:ftruns:ft02}

FT02 continued F001's masked reconstruction on all three sensors - EnMAP, PRISMA and AVIRIS-NG at shares 0.40 /
0.35 / 0.25, one sensor per batch, 20 epochs of 30,000 samples, learning rate $10^{-4}$, 4.25 hours. Its three
pass criteria, written before launch:
\begin{enumerate}
\item Bench mean ROC-AUC (residual Mahalanobis) at least 0.861 at the best checkpoint.
\item No scene below F001's epoch-60 score by more than 0.03.
\item Each sensor's test reconstruction lower at the end than after epoch 1.
\end{enumerate}

\begin{table}[htbp]
\centering
\caption{FT02: validation SAM (degrees) per sensor and the PRISMA bench, ROC-AUC averaged over four scenes.}
\label{tab:ftruns:ft02}
\small
\setlength{\tabcolsep}{4pt}
\begin{tabular}{lccccc}
\toprule
epoch & SAM EnMAP & SAM PRISMA & SAM AVIRIS-NG & res.\ Mahalanobis & combined \\
\midrule
F001 ep109 & -- & -- & -- & 0.862 & 0.706 \\
1 & 4.81 & 7.30 & 7.94 & \textbf{0.864} & 0.760 \\
8 (best.pt) & 4.17 & 5.93 & 6.07 & 0.855 & 0.824 \\
13 & 4.16 & 6.04 & 6.49 & 0.855 & 0.809 \\
20 (last.pt) & 4.07 & 6.31 & 6.19 & 0.853 & 0.814 \\
\bottomrule
\end{tabular}

\smallskip
{\footnotesize Source: \tagM{research/experiments/indradhanu-finetune/FT02-adapt/README.md}.}
\end{table}

\coursefig[0.9\linewidth]{iw-50-ft02-bench}{FT02's bench by epoch against the gate. Residual Mahalanobis (amber)
starts above the 0.861 gate at epoch 1 and settles below it; \texttt{combined} (grey) climbs from F001's 0.706 to
0.824. best.pt is epoch 8, chosen by validation reconstruction, not by the bench.}{fig:ftruns:ft02}

\begin{worked}[label={we:ftruns:gate}]{Reading FT02's gate}
best.pt (epoch 8) scored 0.931, 0.852, 0.757 and 0.879 on the four scenes. Add one at a time:
\begin{align*}
0.931 + 0.852 &= 1.783\\
1.783 + 0.757 &= 2.540\\
2.540 + 0.879 &= 3.419\\
3.419 / 4 &= 0.85475 \approx 0.855
\end{align*}
Criterion 1 asks for 0.861: the miss is $0.861 - 0.855 = 0.006$. Criterion 2 compares each scene with F001 at epoch
60 (0.930, 0.860, 0.775, 0.881):
\begin{align*}
\text{set1:}\ 0.931 - 0.930 &= +0.001 & \text{set2:}\ 0.852 - 0.860 &= -0.008\\
\text{set3:}\ 0.757 - 0.775 &= -0.018 & \text{set4:}\ 0.879 - 0.881 &= -0.002
\end{align*}
The largest drop is 0.018, under 0.03: passed. Criterion 3 (SAM epoch 1 to 20): EnMAP $4.81 - 4.07 = 0.74$, PRISMA
$7.30 - 6.31 = 0.99$, AVIRIS-NG $7.94 - 6.19 = 1.75$ degrees lower: passed. Gate: not met, by 0.006.
\end{worked}

\paragraph{Lesson 1: a gate is read at one checkpoint, chosen by a different rule.}
\index{checkpoint selection}
The bench peaked at epoch 1, at 0.864 - above the gate. It does not count, because the gate is read at the best
checkpoint, and best.pt is chosen by mean validation reconstruction (the fixed ruler of \cref{ch:pretrain}), which
kept improving until epoch 8. Reconstruction and detection pulled apart after the first epoch. Writing ``at the best
checkpoint'' before launch was deliberate: picking whichever epoch scores best on four labelled scenes would tune the
run to those four scenes, turning the bench into training data.

\paragraph{Lesson 2: part of the change is more bands, not better weights.}
The trained-wavelength record of \cref{ch:finetune} grew during FT02 from 202 to 733 wavelengths: AVIRIS-NG added
337, PRISMA 177, EnMAP 17\tagM{research/experiments/indradhanu-finetune/README.md}. Check the arithmetic:
$202 + 337 = 539$, $539 + 177 = 716$, $716 + 17 = 733$. EnMAP added few because SpectralEarth \emph{is} EnMAP, so
most of its wavelengths were already recorded. From epoch 1 on, the bench reads all 185 bands of the PRISMA scenes,
where F001 read 165; any comparison with F001 mixes the two effects.

\paragraph{Lesson 3: \texttt{combined} nearly caught up - and that says something about residual Mahalanobis.}
\texttt{combined} (half L1, half SAM, each divided by the scene's maximum; \cref{ch:residualmaha}) measures how
\emph{big} a pixel's error is. F001 made large ordinary errors on PRISMA's background, which buried the targets.
FT02 trained on PRISMA, PRISMA's validation SAM fell from 7.30 to 5.93 degrees, and the error floor dropped, so
\texttt{combined} rose from 0.706 to 0.824. Residual Mahalanobis had already removed that floor - it subtracts the
local mean residual and whitens by the scene's covariance - so it had little to gain. Keep this in mind: in
\cref{sec:ftruns:absorb} it becomes the central lesson of the whole series.

\paragraph{The see-saw and the decision.} The per-sensor validation did not fall smoothly: in epochs 11--12 and
15--17 PRISMA improved while EnMAP and AVIRIS-NG got worse, then the reverse (epoch 17: PRISMA 5.04, AVIRIS-NG 8.24
degrees). \Cref{ch:finetune} traced it to the model's one BatchNorm and fixed it with mixed batches (FT02b: see-saw
gone, SAM 3.99 / 4.37 / 5.57, bench 0.856, gate still not met). On 2026-10-16 the user decided to \emph{build on}
FT02 rather than discard it: reconstruction was better on every sensor and the miss was small.

\begin{sourcenote}[: the warm-up that does not exist]
FT02's README lists ``LR 1e-4 cosine (1 warm-up epoch)'', and every config sets \texttt{warmup\_epochs: 1}. The
log of 2026-10-16 (\file{research/experiments/indradhanu-finetune/README.md}, FT02b launch entry) records: ``The
configured 1-epoch warm-up does not exist in code (\texttt{warmup\_epochs} is never read).'' The runs had no warm-up.
\end{sourcenote}

\section{Edges: why correcting the score map fails}\label{sec:ftruns:edges}

Look at a residual Mahalanobis score map of a PRISMA bench scene and you see lines: field boundaries, roads,
coastlines. The model reconstructs a sharp boundary slightly blurred or shifted, the error on both sides is large,
and those pixels rank among the top scores. On 2026-10-16 this was measured with rules fixed before the results
(\file{C:/es/edges/edges.py}).

\subsection{Measuring edge strength}

\index{edge strength}\index{gradient!finite difference}
An \emph{edge} is a place where brightness changes quickly from one pixel to the next. For a smooth function $f(x,
y)$, ``how quickly'' is the \emph{gradient}, the vector of partial derivatives $(\partial f / \partial x, \partial f
/ \partial y)$. An image has no derivatives - only values on a grid - so you replace each derivative by a
\emph{finite difference}. Derive the central difference from the definition of the derivative:
\begin{align}
\frac{\partial f}{\partial x}(x, y) &= \lim_{h \to 0} \frac{f(x + h, y) - f(x - h, y)}{2h}.
\intertext{On a pixel grid the smallest step is one pixel, $h = 1$, so drop the limit:}
g_x(i, j) &\approx \frac{f(i, j+1) - f(i, j-1)}{2}, \qquad g_y(i, j) \approx \frac{f(i+1, j) - f(i-1, j)}{2}.
\intertext{The \emph{gradient magnitude} is the length of the vector $(g_x, g_y)$, by Pythagoras:}
\abs{\nabla f}(i, j) &\approx \sqrt{g_x^2 + g_y^2}. \label{eq:ftruns:gradmag}
\end{align}
The diagnosis used a different, related measure: average a pixel's bands into one brightness number, then take the
\emph{largest absolute difference to any valid four-neighbour} (up, down, left, right):
\begin{equation}
e(i, j) = \max\big(\abs{f_{i,j} - f_{i-1,j}},\ \abs{f_{i,j} - f_{i+1,j}},\ \abs{f_{i,j} - f_{i,j-1}},\
\abs{f_{i,j} - f_{i,j+1}}\big). \label{eq:ftruns:edge}
\end{equation}
These are one-sided differences, and taking the maximum instead of a Pythagorean sum is the ``$L_\infty$'' way of
measuring a vector's size. A pixel is \emph{strong} if its $e$ is in the scene's top 10\%.

\begin{worked}[label={we:ftruns:edges}]{Two edge measures on a 3 \texorpdfstring{$\times$}{x} 3 patch}
Band means of a $3 \times 3$ patch with a field boundary running down its right side:
\[
\begin{matrix}
0.10 & 0.10 & 0.30\\
0.10 & 0.12 & 0.30\\
0.10 & 0.10 & 0.31
\end{matrix}
\]
\emph{Central differences at the centre (0.12).} Left 0.10, right 0.30, up 0.10, down 0.10:
\begin{align*}
g_x &= (0.30 - 0.10) / 2 = 0.20 / 2 = 0.10\\
g_y &= (0.10 - 0.10) / 2 = 0\\
\abs{\nabla f} &= \sqrt{0.10^2 + 0^2} = \sqrt{0.01} = 0.10
\end{align*}
\emph{The diagnosis's measure \eqref{eq:ftruns:edge} at the centre.}
\begin{align*}
\abs{0.12 - 0.10} &= 0.02 \quad\text{(up)}\\
\abs{0.12 - 0.10} &= 0.02 \quad\text{(down)}\\
\abs{0.12 - 0.10} &= 0.02 \quad\text{(left)}\\
\abs{0.12 - 0.30} &= 0.18 \quad\text{(right)}\\
e &= \max = 0.18
\end{align*}
The top-left pixel (0.10) has neighbours right 0.10 and down 0.10 inside the patch, so $e = 0$.

\emph{Now a one-pixel target.} Replace the patch by a field of 0.10 with a single bright pixel of 0.30 at the centre.
Central differences at the centre: $g_x = (0.10 - 0.10)/2 = 0$, $g_y = 0$, so $\abs{\nabla f} = 0$ - the
symmetric difference cannot see a one-pixel spike. The one-sided maximum: every neighbour differs by $\abs{0.30 -
0.10} = 0.20$, so $e = 0.20$, a strong edge. (Script check of both\tagC.)
\end{worked}

That last line is the whole trap. A small target, brighter or darker than its field, has a large difference to its
neighbours on every side. By the diagnosis's definition it \emph{is} an edge.

\begin{table}[htbp]
\centering
\caption{Edges and false alarms on the four PRISMA bench scenes (F001 epoch 109 and FT02 epoch 13). By chance each
share would be 10\%, since ``strong'' is the top 10\%.}
\label{tab:ftruns:edges}
\small
\begin{tabular}{lc}
\toprule
measurement & range over the four scenes \\
\midrule
top-1\% false positives on strong edges & 49--68\% \\
ground-truth target pixels on strong edges & 29--79\% \\
detected target pixels (hits) on strong edges & 39--91\% \\
\bottomrule
\end{tabular}

\smallskip
{\footnotesize Source: \tagM{research/experiments/indradhanu-finetune/FT02-adapt/README.md} (Observations, edge diagnosis).}
\end{table}

The first row says edges are the main source of false alarms - five to seven times chance. The second row is the trap.

\subsection{Correcting the score map, and why it collapses}

\index{score-map correction}
The obvious fix: compare each pixel only with pixels of similar edge strength. The measured version sorted the scene
into 20 bins by edge strength and scored each pixel by \emph{its rank minus the median rank of its bin}. A pixel on a
strong edge then has to beat other strong-edge pixels, not the smooth field. It collapsed detection: bench mean 0.86
to 0.67. Work why on eight pixels.

\begin{worked}[label={we:ftruns:bins}]{Rank minus the median rank of its bin}
Eight pixels, two bins. A is a real target on an edge; B, C, D are edge false alarms; E, F, G, H are a smooth field
with small noise. Scores and bins:
\begin{center}
\small\setlength{\tabcolsep}{3pt}
\begin{tabular}{lcccccccc}
\toprule
pixel & A & B & C & D & E & F & G & H \\
\midrule
score & 9 & 10 & 8 & 7 & 3 & 2 & 1 & 4 \\
bin & strong & strong & strong & strong & smooth & smooth & smooth & smooth \\
rank (1 = lowest) & 7 & 8 & 6 & 5 & 3 & 2 & 1 & 4 \\
\bottomrule
\end{tabular}
\end{center}
Medians of the ranks in each bin. Strong: ranks 5, 6, 7, 8; the median of four numbers is the mean of the middle two,
$(6 + 7)/2 = 6.5$. Smooth: ranks 1, 2, 3, 4; median $(2 + 3)/2 = 2.5$.

Corrected scores:
\begin{align*}
A:\ 7 - 6.5 &= 0.5 & B:\ 8 - 6.5 &= 1.5 & C:\ 6 - 6.5 &= -0.5 & D:\ 5 - 6.5 &= -1.5\\
E:\ 3 - 2.5 &= 0.5 & F:\ 2 - 2.5 &= -0.5 & G:\ 1 - 2.5 &= -1.5 & H:\ 4 - 2.5 &= 1.5
\end{align*}
Before, the target A was second of eight. After, B and H share first place (1.5) and A ties with E, a noise pixel,
for third. The correction promoted the top of the smooth field - pixels that were not anomalous at all - to the level
of the strongest edges, and it pushed the target down with its edge neighbours. (Script check\tagC.)
\end{worked}

\begin{table}[htbp]
\centering
\caption{Earlier score-map sharpeners on F001 epoch 52; none beat plain residual Mahalanobis (D).}
\label{tab:ftruns:postscore}
\small
\begin{tabular}{lcc}
\toprule
method, applied to the score map & bench AUC & objects caught of 909 \\
\midrule
\textbf{D} residual Mahalanobis, no post-processing & \textbf{0.865} & 383 \\
F(C) difference-of-Gaussians blob boost & 0.718 & 403 \\
E(C) local z-score of the score map (ring) & 0.790 & -- \\
G-C suppress errors explained by a $\pm 1$\,px shift & 0.845 & -- \\
\bottomrule
\end{tabular}

\smallskip
{\footnotesize Source: \tagM{research/experiments/indradhanu-finetune/FT02-adapt/README.md} (\file{C:/es/postscore/}).}
\end{table}

The blob boost found 20 more objects and gave up $0.865 - 0.718 = 0.147$ of AUC to do it: it lit up many more
non-targets on the way. The FT02 README's conclusion: edge strength, measured on the score map or the image, cannot
tell a field boundary from a target, because at the scale of a few pixels they look the same.

\subsection{An edge term under the blind spot}

\index{blind spot}
Training is a different place to act. Recall inference (\cref{ch:inference}): two complementary checkerboard masks,
so every pixel is reconstructed while its own token - a $4 \times 4$ pixel square, stage 1's patch embedding with
kernel 4 and stride 4 - is hidden. The model predicts the hidden token's pixels from the visible tokens around it.
Compare the two things that look alike on the score map:
\begin{itemize}
\item \textbf{A field boundary} that crosses the hidden token continues on both sides of it. Its direction and its
  contrast are visible in the neighbouring tokens. A model that has learned to continue straight boundaries sharply
  can predict it.
\item \textbf{A small target} sitting inside the hidden token has no continuation. Nothing in the visible tokens says
  it is there. The model predicts the field, the error is large - and that error is the detection.
\end{itemize}
So a loss that rewards \emph{sharp} reconstruction of hidden pixels can teach the first and cannot teach the second.
The score map has already lost this information: it records only how large the error was, not whether the context
could have predicted the structure.

\coursefig[\linewidth]{iw-53-edges}{(a) Edge-strength correction (amber) pushes down everything on strong edges, the
field boundary and the target alike. (b) Under the blind spot: a boundary passing through a hidden token (grey
hatching) can be continued from its visible ends; a target wholly inside the hidden token leaves no trace in the
context.}{fig:ftruns:edges}

\paragraph{The gradient-matching loss.} \index{gradient-matching loss}
FT02c added \code{gradient_l1} (\code{app/foundation_models/trainers/finetune_trainer.py}) to the mixed-batch loss
with weight \texttt{grad\_loss\_weight}. Its docstring: ``mean $\lvert d\,\hat{x} - d\,x \rvert$ over horizontal and
vertical neighbour differences, per band, on pairs whose two pixels are both in loss\_mask \ldots\ and valid in that
band. Blur shows up as a gradient mismatch.'' In symbols, for one band, with $P$ the set of neighbouring pairs
$(p, q)$ both hidden, valid and eroded (\cref{ch:masking}):
\begin{equation}
L_\text{grad} = \frac{1}{\abs{P}} \sum_{(p,q) \in P} \Big\lvert (\hat{x}_q - \hat{x}_p) - (x_q - x_p) \Big\rvert .
\label{eq:ftruns:gradl1}
\end{equation}
L1 on the values (\cref{ch:losses}) punishes blur only mildly; a loss on the differences punishes it directly.

\begin{worked}[label={we:ftruns:gradl1}]{The same blur, two losses}
A row of four hidden pixels across a boundary:
\begin{align*}
\text{input:}\quad & 0.10,\ 0.10,\ 0.50,\ 0.50\\
\text{reconstruction:}\quad & 0.10,\ 0.20,\ 0.40,\ 0.50 \qquad\text{(the boundary, blurred)}
\end{align*}
\emph{Value L1.} Absolute differences: $0,\ 0.10,\ 0.10,\ 0$. Sum $0.20$; mean $0.20 / 4 = 0.05$.

\emph{Gradient L1.} Input differences (next minus this): $0.10 - 0.10 = 0$, $0.50 - 0.10 = 0.40$, $0.50 - 0.50 =
0$. Reconstruction differences: $0.20 - 0.10 = 0.10$, $0.40 - 0.20 = 0.20$, $0.50 - 0.40 = 0.10$. Mismatches:
$\abs{0.10 - 0} = 0.10$, $\abs{0.20 - 0.40} = 0.20$, $\abs{0.10 - 0} = 0.10$. Sum $0.40$; mean $0.40 / 3 = 0.133$.

The same blur costs 0.05 in value L1 and 0.133 in gradient L1: the edge term is what notices that a step became a
ramp. (Script: 0.0500 and 0.1333\tagC.)
\end{worked}

\begin{sourcenote}[: where a hidden target could leak]
It has been said that overlapping patch embeddings leak one or two pixels of a hidden token into its neighbours. In
this model they do not, at the stage that hides tokens: stage 1's embedding is non-overlapping, kernel 4, stride 4,
padding 0 (\code{app/foundation_models/components/seg_former_encoder.py}: ``Stage 1: kernel=4, stride=4
(NON-OVERLAPPING - zero information leakage''), hidden tokens are restored as zeros after stage 1, and the spectral
stem in front works on one pixel at a time. The real leak is a target that spans several tokens: some of them are
visible in each checkerboard pass, so an edge-continuer could ``continue'' it into the hidden part.
\end{sourcenote}

\paragraph{What FT02c showed.} FT02c ran on 2026-10-16 (2.2 hours, code \texttt{9ff0484}) from FT02b's best.pt with
\texttt{grad\_loss\_weight} 1.0 and the SAM weight held at 0.5 throughout. Its criterion: the share of top-1\% false
alarms on strong edges down on at least 3 of 4 scenes.

\begin{table}[htbp]
\centering
\caption{FT02c against its parent FT02b.}
\label{tab:ftruns:ft02c}
\small
\begin{tabular}{lcc}
\toprule
 & FT02b ep9 & FT02c ep10 \\
\midrule
val.\ SAM, EnMAP / PRISMA / AVIRIS-NG & 3.99 / 4.37 / 5.57 & 3.93 / 4.21 / 5.41 \\
edge term \texttt{train\_grad\_l1}, ep 1 $\to$ 10 & -- & 0.01163 $\to$ 0.01153 \\
bench ROC-AUC, residual Mahalanobis & 0.856 & 0.855 \\
pAUC / caught top 1\% / top 0.1\% & 0.158 / 411 / 128 & 0.159 / 413 / 131 \\
top-1\% false alarms on edges, set1--4 & 69 / 52 / 52 / 60\% & 69 / 52 / 52 / 61\% \\
\bottomrule
\end{tabular}

\smallskip
{\footnotesize Source: \tagM{research/experiments/indradhanu-finetune/FT02c-edges/README.md}.}
\end{table}

The edge term barely moved: $0.01163 - 0.01153 = 0.00010$, and $0.00010 / 0.01163 = 0.0086$, about 1\% in ten
epochs, while L1 and SAM kept falling. The edge false-alarm share is unchanged on all four scenes. Verdict: not met.

\paragraph{Why the term could not be reduced.} \index{pixel shuffle}
The decoder (\cref{ch:decoder}) fuses all encoder stages at a quarter of the resolution, $H/4$, and the pixel shuffle
of \citet{shi2016realtime} turns \emph{one} feature vector into a whole $4 \times 4$ block of pixels
(\code{app/foundation_models/components/seg_former_decoder.py}). Where an edge falls inside that block has to be
encoded in that one quarter-resolution vector; with the loss switched on, it was not learned. The likely fix is
architectural - a higher-resolution path into the decoder, masked so it cannot see hidden tokens - not a loss term.
(That path was built later, as FT09; \cref{ch:decoderwork}.) The small reconstruction gains over FT02b are, in the
README's words, ``more plausibly from the constant SAM weight 0.5 (FT02b ramped it from 0) than from the edge term''.

\begin{keypoint}
Edges are the main false alarm and a common shape of a true target; a correction applied to the finished score map
cannot tell them apart, and a loss that could, in principle, was blocked by the decoder's $4 \times 4$ bottleneck.
\end{keypoint}

\section{The options not taken}\label{sec:ftruns:options}

After FT02, four other directions were weighed on 2026-10-16 and set aside, each for a stated reason. Knowing what
each would cost is what lets you pick one up later without repeating the discussion.

\coursefig[\linewidth]{iw-54-options}{The four options discussed after FT02, what each would add, its main cost, and
where it stood.}{fig:ftruns:options}

\paragraph{1. A teacher (``oracle'') head - deferred.} \index{exponential moving average}\index{teacher-student}
A \emph{teacher} is a copy of the model that sees more; a \emph{student} is the model being trained, which sees less
and is trained to match the teacher's output. Here the teacher would read a sensor's full band list and the student a
degraded one. The teacher would be an \emph{exponential moving average} (EMA) of the student's own weights: after
each step, $\theta_\text{T} \leftarrow 0.999\,\theta_\text{T} + 0.001\,\theta_\text{S}$. How slow is that? Hold the
student fixed at $s$ and subtract $s$ from both sides of the update:
\begin{align}
\theta_\text{T}^{(n+1)} - s &= 0.999\,\theta_\text{T}^{(n)} + 0.001\,s - s = 0.999\,\big(\theta_\text{T}^{(n)} - s\big).
\intertext{The gap shrinks by a factor 0.999 per step, so after $n$ steps}
\theta_\text{T}^{(n)} - s &= 0.999^n\,\big(\theta_\text{T}^{(0)} - s\big).
\intertext{The gap halves when $0.999^n = 1/2$. Take logarithms of both sides and divide:}
n &= \frac{\ln 2}{-\ln 0.999} = \frac{0.6931}{0.0010005} \approx 693 \text{ steps.} \label{eq:ftruns:halflife}
\end{align}
One step from teacher 0.50 toward student 0.80: $0.999 \times 0.50 + 0.001 \times 0.80 = 0.4995 + 0.0008 =
0.5003$. The cost: a second forward pass per step, memory for a second copy of the weights, and one more thing that
can drift. It was deferred because a teacher built on FT02's swings would inherit them.

\paragraph{2. Per-sensor heads - not chosen.} A decoder per sensor would give each its own capacity, but it gives up
the design's point: one set of wavelength-aware weights that reconstructs any band list, including sensors nobody has
trained on. It also splits the data, and the effects mixed batches address (BatchNorm, gradient balance) sit
\emph{before} any head.

\paragraph{3. A synthetic-target detection head - not chosen (then).} \index{linear mixture}
Implant artificial targets and train a head to find them. An implanted target is usually a \emph{linear mixture}: at
one band, background 0.20, target 0.50, fraction 0.30 gives $0.70 \times 0.20 + 0.30 \times 0.50 = 0.14 + 0.15 =
0.29$. The cost is in the word ``synthetic'': the head learns whatever the implanting recipe makes detectable - a
perfectly linear mix, a sharp square outline - and it changes what the model is, from ``fails to reconstruct the
unexpected'' to ``knows what a target looks like''.

\paragraph{4. Regularisation - shelved.} \index{L2-SP}
Four standard ways to stop fine-tuning drifting from pretraining: an EMA of the weights, replaying SpectralEarth in
each batch, a lower learning rate, and \emph{L2-SP} (``L2 towards the starting point''), which adds a penalty
\begin{equation}
R(\vect{w}) = \frac{\alpha}{2} \sum_i \big(w_i - w_i^{0}\big)^2 \label{eq:ftruns:l2sp}
\end{equation}
to the loss, with $\vect{w}^0$ the pretrained weights. Differentiate with respect to one weight to see what it does
to training: the power rule brings the 2 down, which cancels the $\tfrac12$,
\[
\frac{\partial R}{\partial w_i} = \alpha\,\big(w_i - w_i^0\big),
\]
a pull back toward the start, proportional to how far the weight has moved. With $\alpha = 0.01$ and two weights moved
from $(0.5, -0.2)$ to $(0.6, -0.1)$: the penalty is $\frac{0.01}{2} \times (0.1^2 + 0.1^2) = 0.005 \times 0.02 =
0.0001$, and each weight is pulled back with gradient $0.01 \times 0.1 = 0.001$. All four answer one question - ``is
the model forgetting?'' - and FT02 did not look like forgetting: reconstruction improved on every sensor, EnMAP
included, the sensor closest to SpectralEarth.

\begin{drift}
Status as of 2026-10-18: option 3 ran, in a different form - not as a head but as implant-denoising (FT05,
\cref{sec:ftruns:ft05}). Option 1's teacher idea ran as a teacher for \emph{masks} (FT04's pseudo-labels,
\cref{sec:ftruns:ft04}); cross-sensor reconstruction distillation is still deferred. Regularisation and per-sensor heads
have no code.
\end{drift}

\section{Stage 3: FT03, masks, and why the implants had to change first}\label{sec:ftruns:stage3}

\paragraph{FT03 passed.} FT03 trained reconstruction and segmentation jointly on mixed batches, from FT02c with FT01's
head and segmentation weight 0.1, 10 epochs, 3.0 hours. All three criteria were met: water / cloud IoU 0.841 / 0.815
(at or above FT01's 0.82 / 0.80), validation SAM within 0.1 degree of FT02c (3.92 / 4.21 / 5.48), and detection
unchanged (PRISMA pAUC 0.158 against FT02c's 0.159)\tagM{research/experiments/indradhanu-finetune/FT03-joint/README.md}.
Joint training changed detection by nothing measurable, as expected with a segmentation weight of 0.1.

\paragraph{The head's first job: exclusion.} Class-aware scoring - a separate ``normal'' per surface class - was tried
and not shipped: it removed cloud false alarms but water's tight statistics made water noise dominate (set1: water's
share of false alarms 14\% to 92\%; PRISMA pAUC 0.158 to 0.134). The user's rule settled it: ``we don't detect
anything on water and cloud.'' So the head marks them and the scorer sets them to 0 after scoring (\code{exclude} in
\code{app/evaluation/scorers.py}). The mask comes from FT03, but it can be applied to \emph{every} model's scores.

\begin{table}[htbp]
\centering
\caption{PRISMA bench, residual Mahalanobis, with and without water + cloud exclusion on FT03's masks.}
\label{tab:ftruns:exclude}
\small
\begin{tabular}{lcccc}
\toprule
model & caught top 0.1\%, plain & excluded & pAUC, plain & excluded \\
\midrule
v0.2.0 & 14.1\% & 15.3\% & 0.159 & 0.165 \\
F001 ep109 & 13.2\% & 14.6\% & 0.148 & 0.155 \\
FT02c ep10 & 14.4\% & 16.1\% & 0.159 & 0.167 \\
FT03 ep10 & 14.3\% & 16.5\% & 0.158 & 0.165 \\
\bottomrule
\end{tabular}

\smallskip
{\footnotesize Source: \tagM{C:/es/eval/stage3/report.md} (the \texttt{exclude(water+cloud;mask)} rows); summarised in \tagM{docs/lld/evaluation.md}.}
\end{table}

The pAUC gains are $0.165 - 0.159 = 0.006$, $0.155 - 0.148 = 0.007$, $0.167 - 0.159 = 0.008$ and $0.165 - 0.158 =
0.007$ - between 0.006 and 0.008 for every model, including v0.2.0, which has no segmentation head and was never
fine-tuned. The gain belongs to the mask, not to any model's reconstruction. Two costs: ROC-AUC falls, and 39 of the
909 real objects lie wholly inside FT03's water / cloud mask. Checked later: 13 of set1's 27 are single pixels more
than 3\,px offshore - likely vessels. The user then put \emph{vessels on water in scope: only cloud is excluded}
(\cref{ch:runeval}). FT03 with cloud only: pAUC 0.158 to 0.166.

\paragraph{The model-free detector, and what was wrong with bench-v1.} The stage-3 run also scored PCA-RX: classical
local RX on the cube's top 30 principal components, no Indradhanu at all. On bench-v1 it caught 30.2\% at the top 1\%
against FT03's 26.9\%; on the real objects its TPR at 0.1\% was 0.029 against FT03's 0.058 - half. bench-v1 rewarded
something PCA-RX has more of than any model. A diagnosis compared the real objects with bench-v1's implants on two
numbers, both starting from $\vect{d}$, an object's mean spectrum minus its ring's:
\begin{itemize}
\item \textbf{Foreign fraction}: the share of $\norm{\vect{d}}^2$ lying outside the scene's 10 main spectral
  directions. Real 0.006, bench-v1 0.042 (easy tier: real 0.002, bench-v1 0.029).
\item \textbf{Halo over object contrast}: the contrast of the one-pixel ring touching the object, over the object's
  own. Real 0.74, bench-v1 0.60 (easy: 0.70 against 0.38). Real objects bleed into their neighbours; implants
  stopped dead at their edge.
\end{itemize}

\coursefig[0.9\linewidth]{iw-57-what-v2-must-match}{Medians on PRISMA, real objects (grey) against bench-v1 implants
(amber), overall and in the easy tier; the shaded band is within 25\% of the real value.}{fig:ftruns:v2match}

\paragraph{Matching foreignness, derived.} \index{foreignness}
bench-v2's fix rescales the part of $\vect{d}$ outside the scene's subspace so that the foreign fraction hits a
target $f$, keeping the inside part (\code{match_foreignness} in \code{app/evaluation/implant.py}). Split $\vect{d} =
\vect{i} + \vect{o}$ into its inside and outside parts; they are perpendicular, so by Pythagoras $\norm{\vect{d}}^2 =
\norm{\vect{i}}^2 + \norm{\vect{o}}^2$. Scale the outside part by $\beta$ and demand the target share:
\begin{align}
\frac{\beta^2 \norm{\vect{o}}^2}{\norm{\vect{i}}^2 + \beta^2\norm{\vect{o}}^2} &= f.
\intertext{Multiply both sides by the denominator:}
\beta^2 \norm{\vect{o}}^2 &= f\,\norm{\vect{i}}^2 + f\,\beta^2 \norm{\vect{o}}^2.
\intertext{Collect the $\beta^2$ terms on the left and factor:}
\beta^2 \norm{\vect{o}}^2 (1 - f) &= f\,\norm{\vect{i}}^2.
\intertext{Divide by $\norm{\vect{o}}^2(1-f)$:}
\beta^2 &= \frac{f}{1 - f} \cdot \frac{\norm{\vect{i}}^2}{\norm{\vect{o}}^2}. \label{eq:ftruns:beta}
\end{align}

\begin{worked}[label={we:ftruns:beta}]{The repo's \texttt{match\_foreignness} test case}
Scene direction = band 0, $\vect{d} = (3, 4, 0)$, target $f = 0.1$. Inside part $\vect{i} = (3, 0, 0)$ with
$\norm{\vect{i}}^2 = 9$; outside part $\vect{o} = (0, 4, 0)$ with $\norm{\vect{o}}^2 = 16$.
\begin{align*}
\beta^2 &= \frac{0.1}{0.9} \times \frac{9}{16} = 0.1111 \times 0.5625 = 0.0625\\
\beta &= \sqrt{0.0625} = 0.25\\
\vect{d}' &= (3,\ 0.25 \times 4,\ 0) = (3, 1, 0)
\end{align*}
Check: outside $1^2 = 1$; total $9 + 1 = 10$; share $1 / 10 = 0.1$. (Script: $\beta = 0.25$\tagC.)
\end{worked}

The acceptance rule, fixed before the build, asked that foreignness and halo medians lie within 25\% of the real
objects', overall and per tier. For the overall foreign fraction that band is $0.006 \times 0.75 = 0.0045$ to $0.006
\times 1.25 = 0.0075$; bench-v1's 0.042 fails. For the easy-tier halo, $0.70 \times 0.75 = 0.525$ to $0.70 \times 1.25
= 0.875$; bench-v1's 0.38 fails, even though its overall halo (0.60, against a band of 0.555 to 0.925) passes. That is
why the test is per tier. bench-v2 was built and accepted the same day; \cref{ch:benchmarks} teaches it.

\begin{keypoint}
No model was trained on detection until the implants looked like real objects: the cheapest way to lower a loss that
rewards finding bench-v1's implants is to become a better PCA-RX, which the real bench says is the wrong direction.
\end{keypoint}

\section{FT05: implant-denoising}\label{sec:ftruns:ft05}

With bench-v2 accepted, option 3 ran in a new form on 2026-10-17: keep the model and the score, put anomalies into the
\emph{input}, and keep the clean tile as the \emph{target}.

\subsection{Why the old target works against detection}

\index{implant-denoising}
Every earlier run used the input as the target: hide most tokens, reconstruct them, compare with what was there.
Anomalies in the training tiles are in the target too. An object smaller than a token is invisible to its own
reconstruction, but an object larger than a token, or one with a blurred halo, leaks into the visible neighbouring
tokens. A model trained on input targets is rewarded for carrying that leak into its prediction.

\begin{worked}[label={we:ftruns:denoise}]{Which prediction does each target reward?}
One band of one pixel. The background under it would be 0.20; an object makes it 0.50; the model sees part of the
object in a neighbouring token and can predict either value. The loss is L1.
\begin{center}
\begin{tabular}{lcc}
\toprule
 & predict 0.20 & predict 0.50 \\
\midrule
input target 0.50 & $\abs{0.20 - 0.50} = 0.30$ & $\abs{0.50 - 0.50} = 0$ $\leftarrow$ rewarded \\
clean target 0.20 & $\abs{0.20 - 0.20} = 0$ $\leftarrow$ rewarded & $\abs{0.50 - 0.20} = 0.30$ \\
score &$\abs{0.50 - 0.20} = 0.30$ (found) & $\abs{0.50 - 0.50} = 0$ (missed) \\
\bottomrule
\end{tabular}
\end{center}
With the input as target, the loss prefers the prediction that hides the object from the score. With the clean tile
as target, it prefers the prediction that exposes it. FT05 changes only the target, and that flips which prediction the
loss rewards.
\end{worked}

\subsection{The implants and the weighted loss}

\code{implant_sample} (\code{app/foundation_models/trainers/finetune_trainer.py}) runs in the data-loader workers on
every \emph{training} tile: with probability 0.5 it writes 1 to 4 bench-v2 implants (\code{implant_tile}, the per-tile
twin of the bench builder), keeps the clean tile as \texttt{clean.npy}, and marks every pixel within 3\,px of an
implant as \emph{near}. The generator is seeded from the tile's key, so a tile gets the same implants every epoch.
Implants come from two sources, \texttt{signature} and \texttt{material}; bench-v2's third source, \texttt{host}, is
held out so that its row of bench-v2 is an out-of-recipe check.

In \code{_mixed_loss} the model reads the implanted tile and is graded against the clean one, and the loss mask is
re-weighted:

\begin{lstlisting}
loss_mask = hidden_pixels * eroded_valid       # (B, 1, H, W), as in every run
# near: (B, 1, H, W), 1 within 3 px of an implant; implant_near_weight = 5
loss_mask = loss_mask * (1.0 + (implant_near_weight - 1.0) * near)
\end{lstlisting}

\paragraph{How many pixels are near?} \index{dilation}
\code{binary_dilation(gt > 0, iterations=3)} grows the implant by the default cross-shaped structuring element three
times. Each pass adds the four direct neighbours of every marked pixel, so after $r$ passes a single pixel has grown
into a diamond: every offset $(\Delta y, \Delta x)$ with $\abs{\Delta y} + \abs{\Delta x} \le r$. Count it ring by
ring: the ring at distance $k \ge 1$ has $4k$ pixels (one per unit of $k$ along each of the four sides), so
\begin{equation}
\abs{\text{diamond}_r} = 1 + \sum_{k=1}^{r} 4k = 1 + 4 \cdot \frac{r(r+1)}{2} = 1 + 2r(r + 1). \label{eq:ftruns:diamond}
\end{equation}
For $r = 3$: $1 + 2 \times 3 \times 4 = 1 + 24 = 25$, which is $1 + 4 + 8 + 12$. SciPy agrees: 25 for a one-pixel
implant and 32 for a two-pixel one\tagC.

\paragraph{Why weight them, and what the weight does.} \index{weighted mean}
\code{_masked_l1_sam} divides by the sum of the mask, so the loss is a \emph{weighted mean} over pixels. If $N$ pixels
count, $n$ of them near an implant with weight $w$ and the rest with weight 1, the near pixels' share of the loss's
weight is
\begin{equation}
\text{share}(w) = \frac{w\,n}{(N - n) + w\,n}. \label{eq:ftruns:share}
\end{equation}

\begin{worked}[label={we:ftruns:nearweight}]{Times five moves 0.15\% to 0.76\%}
A $128 \times 128$ tile has 16,384 pixels; only the hidden ones count (mask ratio 0.65). One one-pixel implant has 25
near pixels.
\begin{align*}
N &= 0.65 \times 16{,}384 = 10{,}650 \quad\text{(hidden pixels)}\\
n &= 0.65 \times 25 = 16.25 \quad\text{(hidden near pixels, on average)}\\
\text{share}(1) &= \frac{16.25}{10{,}650} = 0.0015 \quad (0.15\%)\\
w\,n &= 5 \times 16.25 = 81.25\\
(N - n) + w\,n &= 10{,}650 - 16.25 + 81.25 = 10{,}715\\
\text{share}(5) &= \frac{81.25}{10{,}715} = 0.0076 \quad (0.76\%)
\end{align*}
The weight changes \emph{shares}, not the loss's overall size, because the denominator grows too. The rest of the
tile, and the clean half of the data, keep teaching plain reconstruction. (Script: 0.00153 and 0.00758\tagC.)
\end{worked}

\coursefig[\linewidth]{iw-58-ft05-step}{One FT05 step. The loader workers (CPU) implant half the tiles and keep the clean
copy; the model (GPU) reads the implanted tile and is graded against the clean one, $\times 5$ near the implants. The
CPU side sets the pace.}{fig:ftruns:ft05}

\subsection{The bill: a data loader that could not keep up}

FT03 ran 10 epochs in 3.0 hours, about $3.0 \times 60 / 10 = 18$ minutes an epoch. FT05 ran at about 2 h 40 min =
160 minutes an epoch, with the GPU 0--7\% busy\tagM{research/experiments/indradhanu-finetune/FT05-implant-denoise/README.md}.
The ratio is $160 / 18 = 8.9$; since the GPU's work per epoch did not change, it could be busy at most $18 / 160 =
0.11$ of the time. Ten epochs would take $10 \times 160 = 1{,}600$ minutes $= 26.7$ hours. All the implant work is NumPy
on the CPU, inside the loader workers.

One speed fix went in before the relaunch: the background subspace of each tile was computed by an SVD of up to
100,000 pixels, and commit \texttt{558ff5e} replaced it by the eigenvectors of the bands-by-bands covariance - ``the
same subspace'', 0.25\,s to about 0.02\,s per tile. Why the same? \index{singular value decomposition}

\begin{worked}[label={we:ftruns:svd}]{Right singular vectors are eigenvectors of \texorpdfstring{$\mat{X}\T\mat{X}$}{XtX}}
Let $\mat{X}$ be the centred pixels, $N$ rows (pixels) by $B$ columns (bands), with singular value decomposition
$\mat{X} = \mat{U}\mat{S}\mat{V}\T$ ($\mat{U}\T\mat{U} = \mat{I}$, $\mat{V}\T\mat{V} = \mat{I}$, $\mat{S}$ diagonal).
Form $\mat{X}\T\mat{X}$, using the transpose of a product $(\mat{A}\mat{B})\T = \mat{B}\T\mat{A}\T$:
\begin{align*}
\mat{X}\T\mat{X} &= (\mat{U}\mat{S}\mat{V}\T)\T (\mat{U}\mat{S}\mat{V}\T)
 = \mat{V}\mat{S}\mat{U}\T\mat{U}\mat{S}\mat{V}\T
 = \mat{V}\mat{S}^2\mat{V}\T .
\end{align*}
That is an eigendecomposition: the columns of $\mat{V}$ are eigenvectors of the $B \times B$ matrix $\mat{X}\T\mat{X}$,
with eigenvalues the squared singular values. For a tile, $N = 16{,}384$ and $B = 224$, so the work moves from a
$16{,}384 \times 224$ matrix to a $224 \times 224$ one.

Check on paper with $N = 4$, $B = 2$: rows $(1, 1)$, $(-1, -1)$, $(2, -2)$, $(-2, 2)$ (each column sums to 0, so
centred).
\begin{align*}
(\mat{X}\T\mat{X})_{11} &= 1^2 + (-1)^2 + 2^2 + (-2)^2 = 1 + 1 + 4 + 4 = 10\\
(\mat{X}\T\mat{X})_{12} &= 1 \cdot 1 + (-1)(-1) + 2(-2) + (-2) \cdot 2 = 1 + 1 - 4 - 4 = -6\\
(\mat{X}\T\mat{X})_{22} &= 1 + 1 + 4 + 4 = 10
\end{align*}
$\begin{psmallmatrix}10 & -6\\ -6 & 10\end{psmallmatrix}$ times $(1, -1)$ gives $(16, -16) = 16\,(1, -1)$, and times
$(1, 1)$ gives $(4, 4) = 4\,(1, 1)$: eigenvalues 16 and 4. Their square roots, 4 and 2, are the singular values of
$\mat{X}$ (script check\tagC).
\end{worked}

The loader was still CPU-bound after that fix. The cause, found on 2026-10-18 and recorded in
\code{single_thread_worker}'s docstring: ``18 workers (3 sources $\times$ 6) on 8 vCPUs oversubscribed the CPU about
8-fold; the implant transform ran 4.7\,s per tile instead of 1.8\,s'' - every worker process had its own pool of
maths-library (BLAS) threads. With one thread per worker and 2 workers per source, FT05 resumed at epoch 6 at 43 minutes
an epoch: $161 / 43 = 3.7$ times faster. FT04 inherited the setting.

\subsection{The verdict}

Pass criteria, fixed before launch: (1) the real PRISMA bench decides - not below FT03 on pAUC to 1\%, TPR at 0.1\% and
caught top 0.1\%, and above on at least two of the three; (2) bench-v2 above FT03, with the held-out \texttt{host} row
not falling; (3) validation SAM within 0.1 degree of FT03 (3.92 / 4.21 / 5.48) and water / cloud IoU within 0.02.

\begin{worked}[label={we:ftruns:ft05crit}]{Criterion 1 at epoch 5}
FT03 against FT05 epoch 5 on the real bench (\file{C:/es/eval/ft05_ep5/report.md}):
\begin{align*}
\text{pAUC:}\quad & 0.15801 - 0.15805 = -0.00004 \quad\text{(below, in the fifth decimal)}\\
\text{TPR at 0.1\%:}\quad & 0.0594 - 0.0582 = +0.0012 \quad\text{(above)}\\
\text{caught top 0.1\%:}\quad & 130 - 130 = 0 \quad\text{(equal)}
\end{align*}
Above on 1 of 3: criterion 1 not met. bench-v2 rose (32.3\% to 32.8\%; pAUC 0.164 to 0.167) and the held-out
\texttt{host} row held (29.4\% to 29.6\%) - the pattern the pre-launch note had called ``optimistic''.
\end{worked}

FT05 finished at epoch 10 on 2026-10-18; best.pt stayed at epoch 5 (lowest validation loss), whose detection is the
tie above. Epoch 10 then failed criterion 3: AVIRIS-NG SAM 5.67 against FT03's 5.48, $+0.19$ against a limit of 0.1.
On bench-v2.1 (\cref{ch:sortbug}) epoch 5 caught 35.4\% against FT03's 34.9\% - a tie. Verdict: \textbf{not passed}.

\begin{sourcenote}[: half of FT05's implants were garbled]
Every \texttt{signature} implant FT05 trained on - half of them - carried the descending-wavelength \code{np.interp}
bug of \cref{ch:sortbug}: each arrived as a flat step, though still at real contrast and foreignness. The bug was fixed
(\texttt{b94dbfd}) after FT05's launch. The FT05 README calls a rerun with correct signatures ``the fair test'' and adds
that ``nothing here suggests it would flip the result''. That rerun has not been made.
\end{sourcenote}

\section{The scoreboard}\label{sec:ftruns:board}

Every model on every current benchmark fits in three tables. They read correctly only with chance and the model-free
PCA-RX row beside them.

\begin{table}[htbp]
\centering
\caption{The real PRISMA bench: 909 labelled objects in 4 scenes. Chance: 10.2\% at the top 1\%, 1.1\% at the top
0.1\%. ``Cloud excluded'' zeroes the model's own cloud mask after scoring.}
\label{tab:ftruns:real}
\small\setlength{\tabcolsep}{4pt}
\begin{tabular}{lccccc}
\toprule
model & caught top 1\% & top 0.1\% & pAUC to 1\% & TPR at 0.1\% & ROC-AUC \\
\midrule
v0.2.0 (shipped) & 44.2\% & 14.1\% & 0.159 & 0.058 & 0.854 \\
F001 ep109 & 44.0\% & 13.2\% & 0.148 & 0.056 & 0.861 \\
FT02c ep10 & 45.4\% & 14.4\% & 0.159 & 0.059 & 0.855 \\
FT03 ep10 & 45.4\% & 14.3\% & 0.158 & 0.058 & 0.853 \\
FT05 ep5 & 44.8\% & 14.3\% & 0.158 & 0.059 & 0.853 \\
FT04 ep10 & 45.7\% & 14.9\% & 0.162 & 0.061 & 0.853 \\
FT03, cloud excluded & 46.1\% & 15.8\% & 0.166 & 0.064 & 0.850 \\
\textbf{FT04, cloud excluded} & 46.0\% & 15.8\% & \textbf{0.170} & \textbf{0.065} & 0.852 \\
PCA-RX (no model) & 39.7\% & 8.1\% & 0.128 & 0.029 & 0.839 \\
\bottomrule
\end{tabular}

\smallskip
{\footnotesize Sources: \tagM{C:/es/eval/stage3/report.md}, \tagM{C:/es/eval/ft04/report.md},
\tagM{research/experiments/indradhanu-finetune/FT04-final-joint/README.md}.}
\end{table}

\begin{table}[htbp]
\centering
\caption{Reconstruction (validation SAM, degrees; lower is better) and EnMAP test IoU for the fine-tuning runs.
Earlier models were validated differently and are left out.}
\label{tab:ftruns:recon}
\small
\begin{tabular}{lcccccc}
\toprule
 & \multicolumn{3}{c}{SAM (deg)} & \multicolumn{3}{c}{IoU} \\
run & EnMAP & PRISMA & AVIRIS-NG & water & cloud & shadow \\
\midrule
FT01 (head only) & -- & -- & -- & 0.82 & 0.80 & 0.20 \\
FT02b ep9 & 3.99 & 4.37 & 5.57 & -- & -- & -- \\
FT02c ep10 & 3.93 & 4.21 & 5.41 & -- & -- & -- \\
FT03 ep10 & 3.92 & 4.21 & 5.48 & 0.841 & 0.815 & 0.205 \\
FT05 ep5 (best.pt) & 3.94 & 4.20 & 5.33 & 0.838 & 0.818 & -- \\
FT05 ep10 (last.pt) & 3.90 & 4.20 & \textbf{5.67} & 0.844 & 0.814 & -- \\
\textbf{FT04 ep10} & \textbf{3.88} & \textbf{4.11} & 5.46 & 0.839 & 0.810 & \textbf{0.223} \\
\bottomrule
\end{tabular}

\smallskip
{\footnotesize Source: the run READMEs, \tagM{research/experiments/indradhanu-finetune/}.}
\end{table}

\begin{figure}[htbp]
\centering
\begin{tikzpicture}
\begin{axis}[
  xbar, bar width=5pt, width=0.8\linewidth, height=9.5cm,
  xmin=0, xmax=50, xlabel={objects caught at the top 1\% (\%)},
  ytick={1,...,8}, yticklabels={{PCA-RX},{FT04 + cloud},{FT04},{FT05 ep5},{FT03},{FT02c},{F001},{v0.2.0}},
  ymin=0.4, ymax=8.6,
  legend style={at={(0.5,1.02)},anchor=south,legend columns=2,font=\small},
  nodes near coords, nodes near coords style={font=\tiny, /pgf/number format/fixed, /pgf/number format/precision=1,
    /pgf/number format/fixed zerofill}]
\addplot[fill=maskgrey!70, draw=black!50] coordinates {(44.2,8) (44.0,7) (45.4,6) (45.4,5) (44.8,4) (45.7,3) (46.0,2) (39.7,1)};
\addplot[fill=loss!70, draw=black!50] coordinates {(35.5,8) (33.1,7) (35.1,6) (34.9,5) (35.4,4) (35.1,3) (36.2,2) (33.2,1)};
\draw[dashed, black!60] (axis cs:10.2,0.4) -- (axis cs:10.2,8.6) node[below right,font=\tiny] {chance $\approx 10\%$};
\legend{real PRISMA bench, bench-v2.1}
\end{axis}
\end{tikzpicture}
\caption{Share of objects caught at the top 1\% on the real bench (grey) and bench-v2.1 (amber), drawn to scale from
\cref{tab:ftruns:real} and \cref{tab:sortbug:v21}. Every model is far above chance and above PCA-RX; the models tie each
other.}
\label{fig:ftruns:board}
\end{figure}

\coursefig[0.9\linewidth]{iw-59-scoreboard}{The course's version of the same board: share caught at the top 1\% for
each model and PCA-RX, on the real bench (grey) and bench-v2.1 (amber).}{fig:ftruns:courseboard}

\begin{drift}
These tables are the board of 2026-10-19 (FT04 added after its run on 2026-10-18). Later models - FT06 to FT09 and
Indradhanu Heavy (\cref{ch:uncertainty,ch:decoderwork,ch:heavy}) - are scored on the same benches and are not shown.
\end{drift}

\paragraph{What each step bought, and how big noise is.} Take bench-v2.1's caught-at-top-1\% apart step by step,
with the rough ruler of \cref{ch:sortbug} ($\operatorname{SE} \approx 0.9$ points for $p = 0.35$, $n = 2{,}817$):
\begin{center}
\small
\begin{tabular}{llcl}
\toprule
step & caught & change & reading \\
\midrule
v0.2.0 $\to$ F001 & 35.5 $\to$ 33.1 & $-2.4$ ($2.7$ SE) & pretraining alone: slightly worse \\
F001 $\to$ FT02c & 33.1 $\to$ 35.1 & $+2.0$ ($2.2$ SE) & three real sensors, mixed batches \\
FT02c $\to$ FT03 & 35.1 $\to$ 34.9 & $-0.2$ & joint training: a tie \\
FT03 $\to$ FT04 & 34.9 $\to$ 35.1 & $+0.2$ & pseudo-labels, rare-class weights: a tie \\
FT04 $\to$ + cloud & 35.1 $\to$ 36.2 & $+1.1$ & a mask, not a model \\
FT03 $\to$ FT05 ep5 & 34.9 $\to$ 35.4 & $+0.5$ & implant-denoising: a tie \\
\bottomrule
\end{tabular}
\end{center}
Read the two large steps together: pretraining alone cost a little on both benches (real-bench pAUC 0.159 to 0.148),
and fine-tuning on the real sensors won it back - to where v0.2.0 already was.

\section{FT04: pseudo-labels, rare classes and the release candidate}\label{sec:ftruns:ft04}

FT03's segmentation head learned only from EnMAP's labels. PRISMA and AVIRIS-NG tiles went through the segmentation
pass, but no loss read their outputs. So the head that would mask clouds on a PRISMA scene had never been graded on
one. FT04, the ``final joint'' run, fixed that with pseudo-labels and weighted the rare classes; it ran on 2026-10-18
from FT03's weights for 10 epochs and passed all three of its criteria.

\subsection{Keep only the confident calls}

\index{pseudo-label}
A \emph{pseudo-label} is a label made by a model rather than a person. The teacher is FT03; the student is FT04.
\code{scripts/pseudo_label.py} runs FT03's head over every PRISMA and AVIRIS-NG training tile and stores, per tile, a
\shape{(5, 128, 128)} array of bytes, $\operatorname{round}(255 \times \sigma(z))$ for each class's logit $z$. A
probability of 0.93 is stored as $\operatorname{round}(255 \times 0.93) = \operatorname{round}(237.15) = 237$. The raw
probabilities are kept on purpose, so a threshold can change without re-running the teacher.

\code{pseudo_targets} (\code{app/utils/segmentation/targets.py}) turns them into \emph{targets} (0 or 1 per class and
pixel) and \emph{known} (whether that pixel enters the loss):

\begin{lstlisting}
PSEUDO_CONF = 0.9              # counts only where p > 0.9 or p < 0.1
p = probs.astype(np.float32) / 255.0                       # (5, 128, 128)
targets = (p > conf).astype(np.float32)
known = ((p > conf) | (p < 1 - conf)) & valid[None]   # valid: enough finite bands
known[2:] = False                         # thin cloud, shadow, snow: not taught
not_water = (targets[0] > 0) & (band_mean(refl, wl, 800, 900) > 0.15)
targets[0][not_water] = 0.0               # still known: a known NOT-water
\end{lstlisting}

In stored units, $0.9 \times 255 = 229.5$ and $0.1 \times 255 = 25.5$: ``yes'' needs a byte of 230 or more, ``no'' a
byte of 25 or less. Only water and cloud are taught; the other classes stay unknown on PRISMA and AVIRIS-NG because
FT01 to FT03 never got them good enough to copy.

\paragraph{The water guard.} \index{water guard}
Water is dark in the near-infrared (NIR). A pixel the teacher calls water with a mean 800--900\,nm reflectance above
0.15 is land-bright, so it becomes a \emph{known not-water} - taught ``no'', not left out. The 0.15 is looser than the
implant builder's water rule ($\text{NIR} < 0.08$): it vetoes only calls that are far too bright for water. The source
records the value (the config's description: ``water guarded by NIR $\le$ 0.15'') but not how it was chosen.

\begin{worked}[label={we:ftruns:pseudo}]{Four pixels through \texttt{pseudo\_targets}}
Stored water bytes and NIR means for a $2 \times 2$ tile. Divide each byte by 255, compare with 0.9 and 0.1, then
apply the guard.
\begin{center}
\small
\begin{tabular}{lcccl}
\toprule
pixel & byte & $p = $ byte$/255$ & NIR & decision \\
\midrule
A & 240 & 0.941 & 0.04 & $> 0.9$, NIR $\le 0.15$ $\to$ target 1, known \\
B & 128 & 0.502 & 0.05 & between 0.1 and 0.9 $\to$ unknown \\
C & 13 & 0.051 & 0.30 & $< 0.1$ $\to$ target 0, known \\
D & 235 & 0.922 & 0.22 & $> 0.9$ but NIR $> 0.15$ $\to$ target 0, known (not-water) \\
\bottomrule
\end{tabular}
\end{center}
Three of the four pixels teach; one says ``water'', two say ``not water''; one is left out. Pixel D is the guard at
work. A check before launch, on 150 EnMAP tiles against EnMAP's own labels, gave the guarded pseudo-label water IoU
0.875, and the teacher's raw confident cloud 0.82\tagM{research/experiments/indradhanu-finetune/FT04-final-joint/README.md}.
\end{worked}

\paragraph{Why there is no cloud guard.} A cloud guard, \code{cloud_plausible}, was built and dropped. It required a
mean 450--650\,nm reflectance above 0.2 \emph{and} spectral flatness: the minimum of four window means (460--500,
540--580, 640--680, 840--880\,nm) over their maximum at least 0.7. It rejected 69\% of EnMAP's labelled clouds on a test
shard, and the error it targeted - ``bright desert called cloud'' on set4 - was not there: set4's ``11\% cloud'' is a
cloud / smoke streak and haze. Haze is not flat: with window means 0.30, 0.27, 0.24, 0.18 (falling, because haze
scatters blue most), $\min/\max = 0.18 / 0.30 = 0.60 < 0.7$, so it is rejected, though its visible mean $(0.30 + 0.27 +
0.24)/3 = 0.27$ is bright enough. A guard that throws out most of the true class costs more than the error it blocks.

\subsection{Rare classes, weighted three times}

\index{binary cross-entropy!weighted}
\texttt{seg\_class\_weights} is $[1, 1, 1, 3, 3]$ for water, cloud, thin cloud, shadow and snow. Recall the loss of
\cref{ch:seghead}. For one class with predicted probability $p$ and target $y \in \{0, 1\}$, the binary cross-entropy is
the negative log of the probability the model gave to the right answer:
\begin{equation}
\ell(y, p) = -\big(y \ln p + (1 - y)\ln(1 - p)\big). \label{eq:ftruns:bce}
\end{equation}
(If $y = 1$ the second term vanishes and $\ell = -\ln p$; if $y = 0$ the first vanishes and $\ell = -\ln(1-p)$.)
\code{masked_bce_dice} averages each class's BCE over its known pixels, adds a soft Dice term, and then combines the
classes with a weighted mean over the classes $H$ that have any known pixel in the batch:
\begin{equation}
L = \frac{\sum_{c \in H} w_c\,L_c}{\sum_{c \in H} w_c} . \label{eq:ftruns:wmean}
\end{equation}
In the code: \code{return (w * per)[has].sum() / w[has].sum()}.

\begin{worked}[label={we:ftruns:wbce}]{Weighted BCE on a batch with no snow}
To keep the arithmetic short, use the BCE part only (the Dice term is added to each $L_c$ in the same way and is
weighted identically). Each class has two known pixels; snow has none.
\begin{align*}
\text{water:}\quad & y = 1, p = 0.9: -\ln 0.9 = 0.1054; \quad y = 0, p = 0.2: -\ln 0.8 = 0.2231;\\
 & L_\text{water} = (0.1054 + 0.2231)/2 = 0.1643\\
\text{cloud:}\quad & y = 1, p = 0.7: -\ln 0.7 = 0.3567; \quad y = 0, p = 0.1: -\ln 0.9 = 0.1054;\\
 & L_\text{cloud} = (0.3567 + 0.1054)/2 = 0.2310\\
\text{thin cloud:}\quad & y = 0, p = 0.4: -\ln 0.6 = 0.5108; \quad y = 1, p = 0.6: -\ln 0.6 = 0.5108;\\
 & L_\text{thin} = 0.5108\\
\text{shadow:}\quad & y = 1, p = 0.3: -\ln 0.3 = 1.2040; \quad y = 0, p = 0.2: -\ln 0.8 = 0.2231;\\
 & L_\text{shadow} = (1.2040 + 0.2231)/2 = 0.7136
\end{align*}
Snow has no known pixel, so it drops out of numerator and denominator.
\begin{align*}
\text{unweighted:}\quad & (0.1643 + 0.2310 + 0.5108 + 0.7136) / 4 = 1.6197 / 4 = 0.4049\\
\text{weighted:}\quad & (0.1643 + 0.2310 + 0.5108 + 3 \times 0.7136) / (1 + 1 + 1 + 3)\\
 &= (0.9061 + 2.1408) / 6 = 3.0469 / 6 = 0.5078
\end{align*}
(Script: 0.4049 and 0.5078\tagC.) The total moved only from 0.40 to 0.51. What moved a lot is where the gradient goes:
differentiate \eqref{eq:ftruns:wmean} with respect to one class's loss,
\[
\frac{\partial L}{\partial L_c} = \frac{w_c}{\sum_{h \in H} w_h},
\]
so shadow's share of every gradient step is $1/4 = 0.25$ unweighted and $3/6 = 0.50$ weighted - doubled.
\end{worked}

On a PRISMA or AVIRIS-NG sub-batch only water and cloud are known, both weight 1, so the weights act on EnMAP's batches
only. Result: shadow IoU 0.205 to 0.223.

\subsection{Results, and the 0.159 to 0.170 decomposition}

\begin{worked}[label={we:ftruns:ft04crit}]{FT04's criteria}
\emph{1. Masks: EnMAP IoU not below FT03 by more than 0.01.} Water $0.839 - 0.841 = -0.002$; cloud $0.810 - 0.815 =
-0.005$. Both within 0.01: met.

\emph{3. Reconstruction: SAM within 0.1 degree of FT03.} EnMAP $3.88 - 3.92 = -0.04$; PRISMA $4.11 - 4.21 = -0.10$;
AVIRIS-NG $5.46 - 5.48 = -0.02$. All better: met.

\emph{2. Detection: default not below FT03, and cloud exclusion on FT04's masks at least as good as on FT03's.}
Default: pAUC 0.162 against 0.158, TPR at 0.1\% 0.061 against 0.058, caught top 0.1\% 14.9\% against 14.3\% - at or
above on all three. Exclusion: 0.170 with FT04's masks against 0.166 with FT03's. Met.
\end{worked}

The run also swapped the masks between the two models, which turns two numbers into a small experiment. A
\emph{two-by-two factorial} design varies two factors - here the reconstruction model and the mask source - in every
combination, so each factor's effect can be read with the other held fixed. \index{factorial design}

\begin{table}[htbp]
\centering
\caption{Real-bench pAUC to 1\% with cloud exclusion: every combination of model and mask.}
\label{tab:ftruns:swap}
\small
\begin{tabular}{lcc}
\toprule
 & FT03's masks & FT04's masks \\
\midrule
FT03's reconstruction & 0.166 & 0.165 \\
FT04's reconstruction & 0.170 & 0.170 \\
\bottomrule
\end{tabular}

\smallskip
{\footnotesize Source: \tagM{C:/es/eval/ft04/report.md}.}
\end{table}

\begin{worked}[label={we:ftruns:factorial}]{Main effects of model and mask}
The \emph{main effect} of a factor is the average over the other factor of the change it causes.

Model (FT03 to FT04), averaged over both masks:
\[
\frac{0.170 + 0.170}{2} - \frac{0.166 + 0.165}{2} = 0.170 - 0.1655 = +0.0045 .
\]
Masks (FT03's to FT04's), averaged over both models:
\[
\frac{0.165 + 0.170}{2} - \frac{0.166 + 0.170}{2} = 0.1675 - 0.168 = -0.0005 .
\]
The masks are interchangeable (FT04's exclude less: ROC-AUC 0.852 against 0.850, at no cost at the top). The step under
exclusion is FT04's reconstruction, not its masks. (Script: $+0.0045$ and $-0.0005$\tagC.)
\end{worked}

Now take apart the real-bench gain of the release candidate over what shipped:
\begin{align*}
\text{shipped v0.2.0, default} &\quad 0.159\\
\text{FT04, default} &\quad 0.162 \qquad \text{model: } 0.162 - 0.159 = +0.003\\
\text{FT04 + cloud exclusion} &\quad 0.170 \qquad \text{exclusion: } 0.170 - 0.162 = +0.008\\
\text{total} &\quad +0.011 \qquad \text{exclusion share } 0.008 / 0.011 = 0.73
\end{align*}

\coursefig[\linewidth]{iw-60-ft04}{(a) How one teacher probability becomes a target: taught ``yes'' above 0.9, ``no''
below 0.1, nothing between, and a confident ``water'' over a bright NIR turned into a known not-water. (b) The real-bench
pAUC from the shipped v0.2.0 to the release candidate: the model adds 0.003, cloud exclusion 0.008.}{fig:ftruns:ft04}

And the comparison is not like for like: v0.2.0 was scored without exclusion. The like-for-like run was made on
2026-10-19 (\file{configs/eval/release_like_for_like.json}): both with cloud exclusion on FT04's masks, v0.2.0 scored
0.165 and FT04 0.170 on the real bench; on bench-v2.1 v0.2.0 + exclusion is slightly ahead (36.7\% against
36.2\%)\tagM{research/experiments/indradhanu-finetune/FT04-final-joint/README.md}. A detection tie.

\subsection{Why better reconstruction does not become better detection}\label{sec:ftruns:absorb}

\index{residual Mahalanobis!scale invariance}
The log's conclusion of 2026-10-18: ``Residual Mahalanobis absorbs what better reconstruction changes.'' Here is the
mechanism, as mathematics. Residual Mahalanobis scores a pixel's residual $\vect{r}$ by its squared Mahalanobis distance
in units of the background residual's covariance $\Cov$ (\cref{ch:residualmaha}), $s = \vect{r}\T\Cov\inv\vect{r}$
(after subtracting the local mean). Suppose a better model shrinks \emph{every} residual by the same factor $c$, so
$\vect{r} \to c\,\vect{r}$. The covariance is fitted on those residuals, and scaling every vector by $c$ scales every
product of two components by $c^2$, so $\Cov \to c^2\Cov$ and $\Cov\inv \to c^{-2}\Cov\inv$. Then
\begin{align*}
s' &= (c\,\vect{r})\T (c^2\Cov)\inv (c\,\vect{r})\\
 &= c \cdot c^{-2} \cdot c\ \vect{r}\T\Cov\inv\vect{r}\\
 &= \vect{r}\T\Cov\inv\vect{r} = s .
\end{align*}
The score does not move. In one dimension, with a background spread of 0.010 and a target residual of 0.050 the target
sits $0.050 / 0.010 = 5.0$ spreads out; a better model with spread 0.008 and target residual 0.040 gives $0.040 /
0.008 = 5.0$ spreads - the same\tagA. Only a model that improves on the background \emph{and not} on the targets raises
the score. That is exactly what FT05 tried to teach, without effect.

\begin{keypoint}
What the pretraining and FT01 to FT05 bought: masks (cloud exclusion, about a point for any model), native bands,
better reconstruction and a benchmark you can trust - not detection by the model alone.
\end{keypoint}

FT04 + cloud exclusion on its own masks became the release candidate, and shipped as Indradhanu V2 - the next chapter.

\section{What the runs taught}\label{sec:ftruns:lessons}

\begin{enumerate}
\item \textbf{A class that stays flat while all others climb is a label ceiling} (FT01, thin cloud).
\item \textbf{Read a gate where you said you would}, and expect reconstruction and detection to part company (FT02).
\item \textbf{Count bands before crediting weights}: a growing trained-wavelength record changes what the model reads
  (FT02).
\item \textbf{The score map cannot tell a boundary from a target}; a loss under the blind spot could, but the decoder's
  $4 \times 4$ bottleneck blocked it (FT02c).
\item \textbf{A shelved option is shelved for a reason}; write the reason down so the next person need not re-argue it.
\item \textbf{Fix the benchmark before you train on it} (stage 3, bench-v1 against PCA-RX).
\item \textbf{The target decides what the loss rewards} - and a weight changes shares, not totals (FT05).
\item \textbf{The CPU can be the bottleneck of a GPU run} (FT05's loader).
\item \textbf{A mask can lift every model alike; credit the mask, not the model} (exclusion; FT04's swap).
\item \textbf{A scale-invariant score absorbs uniform improvements} (residual Mahalanobis).
\end{enumerate}

\begin{confusions}
\confusion{FT02 forgot how to detect.}{Not shown. Reconstruction improved on every sensor, and the bench moved 0.862 to
0.855, smaller than some single-scene swings. ``Forgetting'' would predict worse reconstruction on what F001 knew; EnMAP,
the closest sensor to SpectralEarth, got better.}
\confusion{Epoch 1 passed the gate, so FT02 passed.}{No. The gate names the checkpoint: the best by validation
reconstruction. Choosing a checkpoint by the bench would turn the bench into training data.}
\confusion{Edges are false alarms, so suppress edges.}{Edges are the main false alarm and a common shape of a true
target. Suppressing them costs both (0.86 to 0.67).}
\confusion{With a clean target the model never sees anomalies.}{It sees them in every implanted input. It is graded on
what was underneath.}
\confusion{The x5 near weight makes the loss five times bigger.}{The weighted mean divides by the weighted count. x5
shifts the near pixels' share from 0.15\% to 0.76\%; the total stays a mean.}
\confusion{Unknown means not-water.}{No. Unknown pixels never enter the loss. A known not-water (pixel D) is taught as 0.}
\confusion{The 3x class weight makes the loss three times bigger.}{It is a weighted mean: the denominator grows too (6
instead of 4). It doubles shadow's share of the gradient.}
\confusion{Cloud exclusion is FT04's improvement.}{The mask is FT04's head, but the gain is the mask's: it lifts every
model alike, v0.2.0 included, and FT03 scores the same with either head's masks.}
\confusion{FT04 is the best detector.}{On the real bench by 0.003 pAUC without exclusion, and like for like with
exclusion it ties v0.2.0. It is the release candidate because its masks, reconstruction and native bands are the best.}
\end{confusions}

\begin{checkyourself}
\item\label{ex:ftruns:1} The segmentation head of a smaller model has encoder stages of 16, 32, 80 and 128 channels,
  decoder width 64 and 5 classes, with the same layers as \cref{we:ftruns:seghead}. Count its parameters by hand.
\item\label{ex:ftruns:2} Compute the edge strength \eqref{eq:ftruns:edge} and the central-difference gradient
  magnitude \eqref{eq:ftruns:gradmag} of the right-hand centre pixel (0.30) of the patch in \cref{we:ftruns:edges}. Its
  neighbours inside the patch are up 0.30, down 0.31, left 0.12 (treat the missing right neighbour as equal to the pixel
  itself for the central difference).
\item\label{ex:ftruns:3} Redo \cref{we:ftruns:gradl1} for a reconstruction $0.10, 0.10, 0.45, 0.50$. Give the value L1
  mean and the gradient L1 mean.
\item\label{ex:ftruns:4} A two-pixel implant has 32 pixels within 3\,px. With mask ratio 0.65, 10,650 hidden pixels and
  weight 5, what share of a tile's loss weight do its near pixels carry? What weight would give them a 2\% share?
\item\label{ex:ftruns:5} Stored water bytes 199, 231, 24 and 26 for four pixels with NIR 0.05, 0.18, 0.30, 0.04. Give
  each pixel's water target and known flag.
\item\label{ex:ftruns:6} Per-class losses water 0.2, cloud 0.2, thin cloud 0.4, shadow 0.6, no known snow. Compute the
  unweighted and the $[1, 1, 1, 3, 3]$-weighted loss, and shadow's share of the gradient in each case.
\item\label{ex:ftruns:7} Show that the teacher of \cref{sec:ftruns:options} closes 90\% of its gap to a fixed student
  in about 2,300 steps.
\end{checkyourself}
\endgroup

\nextchapter{ch:shipping}
